%CAPA
\begin{titlepage}

%Filia��o
%--------------------------------------------------------------------
\begin{figure}[!ht]
\begin{minipage}[b]{0.7\linewidth}
\begin{tiny}
Laborat�rio de Modelagem, An�lise e Controle de Sistemas N�o-Lineares \\
Departamento de Engenharia Eletr�nica\\
Universidade Federal de Minas Gerais\\
Av. Ant�nio Carlos 6627, 31270-901 Belo Horizonte, MG Brasil\\
Fone: +55 3499-4866 - Fax: +55 3499-4850 %\\
%emmendes@cpdeee.ufmg.br
\end{tiny}
\end{minipage}\hfill
\begin{minipage}[c]{0.25\linewidth}
\begin{flushright}
\vspace{-7.5cm}
\includegraphics[width=1.75\linewidth]{figuras/mac.eps}
\end{flushright}
\end{minipage}
\end{figure}
%--------------------------------------------------------------------

%T�tulo
%--------------------------------------------------------------------
\vspace{1.75cm}
\begin{center}
\thickhrulefill
\par\nobreak
\vspace*{10\p@}%
\hrule
\vspace*{10\p@}%
{\bf \Huge \bfseries Respostas em frequ�ncia generalizadas: Implica��es na est�tica, na din�mica e na an�lise de sistemas n�o-lineares}
\vspace*{10\p@}%
\hrule
\vspace*{40\p@}%
%--------------------------------------------------------------------

%Autor
%--------------------------------------------------------------------
{\large \bf Tales Argolo Jesus}\\[0.5cm]
\end{center}
%--------------------------------------------------------------------


%Descri��o
%--------------------------------------------------------------------
\vspace{1.75cm}
\begin{flushright}
\begin{minipage}{0.7\linewidth}
Tese submetida � banca examinadora designada pelo Colegiado do
Programa de P�s-Gradua��o em Engenharia El�trica da Universidade Federal
de Minas Gerais, como parte dos requisitos necess�rios � obten��o do grau
de Doutor em Engenharia El�trica.\\
\end{minipage}
\end{flushright}
\vspace{1.6cm}


%--------------------------------------------------------------------

%Orienta��o
%--------------------------------------------------------------------
\begin{tabular}{ll}
  {\bf Orientador:} & Prof. Eduardo Mazoni Andrade Mar�al Mendes, PhD \\
\end{tabular}
\vspace{1.6cm}

%--------------------------------------------------------------------
%Cidade e Data
%--------------------------------------------------------------------
\centerline{Belo Horizonte, outubro de 2012}
%--------------------------------------------------------------------
\end{titlepage}
